\section{SHRINCS}
This section specifies the complete SHRINCS signature scheme by composing the building blocks defined in previous sections.

\subsection{Key Generation}
The \texttt{SHRINCS.KeyGen} function generates a fresh key pair and initializes the signing state. It samples a 3n-byte random seed, partitions it into \texttt{SK.seed} (for key derivation), \texttt{SK.prf} (for message randomization), and \texttt{PK.seed} (for domain separation). It then computes \texttt{PK.sf} (the stateful UXMSS root) and \texttt{PK.sl} (the stateless hypertree root), combining them into \texttt{PK.root} via a tweakable hash. The initial state sets \texttt{q=0} with \texttt{valid=true}, indicating the stateful path is available starting from signature number 1.

\begin{verbatim}
    Algorithm: SHRINCS.KeyGen()
    Output:
        SK: secret key
        PK: public key
        state: initial signing state

    1. seed ←$ {0,1}^(3n)
    2. SK.seed ← seed[0:n]
    3. SK.prf ← seed[n:2n]
    4. PK.seed ← seed[2n:3n]
    5. ADRS ← new_ADRS()
    6. PK.sf ← uxmss_root(SK.seed, PK.seed, ADRS)
    7. setLayerAddress(ADRS, d - 1)
    8. setTreeAddress(ADRS, 0)
    9.  PK.sl ← xmss_root(SK.seed, PK.seed, ADRS, h_prime)
    10. setLayerAddress(ADRS, 0)
    11. setTreeAddress(ADRS, 0)
    12. setTypeAndClear(ADRS, ROOT)
    13. PK.root ← Tw(PK.seed, ADRS, PK.sf || PK.sl)
    14. SK ← (SK.seed, SK.prf, PK.seed, PK.sf, PK.sl, PK.root)
    15. PK ← (PK.seed, PK.root)
    16. state ← { q: 0, valid: true }
    17. return (SK, PK, state)
\end{verbatim}

\subsection{Recovery from Seed}
The \texttt{SHRINCS.Restore} function recovers a key pair from a backed-up seed. This enables the standard Bitcoin workflow of restoring a wallet from a seed phrase. However, since the signing state (which \texttt{q} values have been used) cannot be recovered from the seed alone, the stateful path is disabled by setting \texttt{valid=false}. The recovered wallet can only use stateless signatures, which are larger but do not require state tracking. This is the key tradeoff that enables static backups while maintaining security.
\begin{verbatim}
    Algorithm: SHRINCS.Restore(seed)
    Input:
        seed: master seed
    Output:
        SK: secret key
        PK: public key
        state: invalid state (stateless only)

    1. SK.seed ← seed[0:n]
    2. SK.prf ← seed[n:2n]
    3. PK.seed ← seed[2n:3n]
    4. ADRS ← new_ADRS()
    5. PK.sf ← uxmss_root(SK.seed, PK.seed, ADRS)
    6. setLayerAddress(ADRS, d - 1)
    7. setTreeAddress(ADRS, 0)
    8.  PK.sl ← xmss_root(SK.seed, PK.seed, ADRS, h_prime)
    9.  setLayerAddress(ADRS, 0)
    10. setTreeAddress(ADRS, 0)
    11. setTypeAndClear(ADRS, ROOT)
    12. PK.root ← Tw(PK.seed, ADRS, PK.sf || PK.sl)
    13. SK ← (SK.seed, SK.prf, PK.seed, PK.sf, PK.sl, PK.root)
    14. PK ← (PK.seed, PK.root)
    15. state ← { q: false, valid: false }
    16. return (SK, PK, state)
\end{verbatim}

\subsection{Stateful signing}
The \texttt{SHRINCS.SignStateful} function signs a message using the efficient stateful (UXMSS) path. It first checks that state is valid and that signature slots remain (\texttt{q <= hsf + 1}). It then increments \texttt{q} and updates the state before generating the signature. This ordering is critical: if the process is interrupted after signing but before state persistence, the one-time key would be consumed without the state reflecting it, risking catastrophic key reuse on retry (see Section~1.2).

The signature includes \texttt{PK.sl} prepended to enable unified verification. This produces minimal signatures (324 bytes for q = 1 in SHRINCS-B) that grow linearly with \texttt{q}. If state is invalid or exhausted, the function returns "false" and the caller should use stateless signing instead.
\begin{verbatim}
    Algorithm: SHRINCS.SignStateful(m, SK, state)
    Input:
        m: message to sign
        SK: secret key (SK.seed, SK.prf, PK.seed, PK.sf, PK.sl, PK.root)
        state: current signing state { q, valid }
    Output:
        sig: stateful signature, or $\bot$ if state is invalid/exhausted
        state_prime: updated state

    1. if not state.valid:
        return (false, state)
    2. q ← state.q + 1
    3. if q > hsf + 1:
        return (false, state)
    4. state_prime ← { q: q, valid: true }
    5. ADRS ← new_ADRS()
    6. uxmss_sig ← uxmss_sign(m, SK.seed, SK.prf, PK.seed, PK.root, ADRS, q)
    7. sig ← PK.sl || uxmss_sig
    8. return (sig, state_prime)
\end{verbatim}

\subsection{Stateless signing}
The \texttt{SHRINCS.SignStateless} function signs a message using the stateless (SPHINCS+C) path. This is used when state is lost, corrupted, or exhausted.

SHRINCS uses a unified digest for the stateless path: a single call to \texttt{H\_msg\_fors} produces a digest that encodes both the FORS leaf indices (determining which leaves to reveal in the few-time signature) and the hypertree tree/leaf indices (determining which XMSS trees and WOTS+C key pairs to use at each layer). The grinding counter \texttt{ctr} is iterated until the grinding condition is satisfied on this unified digest.

The digest output length is \texttt{(roundup(k*a + hsl) / 8)} bytes. The first \texttt{k*a} bits encode the \texttt{k} FORS leaf indices (with the last a bits required to be zero by grinding). The remaining \texttt{hsl} bits encode the hypertree path indices.

The signature is constructed bottom-up: first a FORS+C signature on the message, then \texttt{d} XMSS signatures on successive layer roots ascending through the hypertree. The signature includes \texttt{PK.sf} prepended to enable unified verification.

\begin{verbatim}
    Algorithm: SHRINCS.SignStateless(m, SK)
    Input:
        m: message to sign
        SK: secret key
    Output:
        sig: stateless signature

    1.  ADRS ← new_ADRS()
    2.  setLayerAddress(ADRS, 0)
    3.  setTreeAddress(ADRS, 0)
    4.  (fors_sig, digest) ← fors_sign(m, SK.seed, SK.prf, PK.seed, PK.root, ADRS)
    5.  (tree_idx, leaf_idx) ← parse_idx(digest)
    6.  indices ← fors_msg_to_indices(digest)
    7.  setLayerAddress(ADRS, 0)
    8.  setTreeAddress(ADRS, tree_idx[0] * 2^h_prime + leaf_idx[0])
    9.  fors_pk ← fors_pkFromSig(fors_sig, indices, PK.seed, ADRS)
    10. ht_sig ← []
    11. msg ← fors_pk
    12. for layer from 0 to d - 1:
        a. setLayerAddress(ADRS, layer)
        b. setTreeAddress(ADRS, tree_idx[layer])
        c. xmss_sig ← xmss_sign(msg, SK.seed, SK.prf, PK.seed, PK.root, ADRS, h', leaf_idx[layer])
        d. ht_sig ← ht_sig || xmss_sig
        e. if layer < d - 1:
               msg ← xmss_root(SK.seed, PK.seed, ADRS, h') // we may cache or reuse the root computed during xmss_sign 
    13. sig ← PK.sf || fors_sig || ht_sig
    14. return sig
\end{verbatim}

\subsection{Index Parsing}
The \texttt{parse\_idx} function extracts tree and leaf indices for each hypertree layer from the stateless digest. In the unified digest scheme, these bits follow the \texttt{k*a} bits used for FORS indices. The \texttt{hsl} bits are partitioned into \texttt{d} pairs of \texttt{(leaf\_idx, tree\_idx)}, extracted from least significant (layer 0) to most significant (layer \texttt{d-1}), with \texttt{h'} bits per leaf index.
%\mknote{This is a bit confusin, that we extract the indices of the OTS/FTS schemes used, but not the indices of the revealed leaves in FORS. With the Hashing changes, it needs a double-check that everything is computed correctly and the same data is not used for two purposes. }
\begin{verbatim}
    Algorithm: parse_idx(digest)
    Input:
        digest: message digest
    Output:
        tree_idx: array of d tree indices
        leaf_idx: array of d leaf indices

    1. ht_offset ← k * a
    2. ht_bits ← extract_bits(digest, ht_offset, hsl)
    3. idx ← ht_bits
    4. tree_idx ← []
    5. leaf_idx ← []
    6. for layer from 0 to d - 1:
        a. leaf_idx[layer] ← idx AND (2^h_prime - 1)
        b. idx ← idx >> h_prime
        c. tree_idx[layer] ← idx
    7. return (tree_idx, leaf_idx)
\end{verbatim}
\paragraph{Note:} The total digest must encode \texttt{k*a + hsl = 6*22 + 24 = 156} bits. The \texttt{H\_msg\_fors} output length (Section~5.4.2) must be set accordingly: \texttt{roundup((k*a + hsl) / 8) = 20} bytes. This is a correction from the original \texttt{roundup((k*a) / 8) = 17} bytes, which was insufficient for the unified digest.

\subsection{Stateful verification}
The \texttt{SHRINCS.VerifyStateful} function verifies a stateful signature. It extracts \texttt{PK.sl} from the signature, parses the UXMSS signature, and infers \texttt{q} from the authentication path length. It then uses the two-phase WOTS+C hashing to recover \texttt{PK.sf} via \texttt{uxmss\_pkFromSig} and verifies that \texttt{Tw(PK.seed, ADRS, PK.sf || PK.sl)} equals \texttt{PK.root}.
\begin{verbatim}
    Algorithm: SHRINCS.VerifyStateful(m, sig, PK)
    Input:
        m: message
        sig: stateful signature
        PK: public key
    Output:
        valid: boolean

    1. ADRS ← new_ADRS()
    2. PK.sl ← sig[0 : n]
    3. uxmss_sig ← sig[n : ]
    4. wots_sig_size ← R_SIZE + c + l * n
    5. auth_len ← len(uxmss_sig) - wots_sig_size
    6. if auth_len < 0 or (auth_len mod n) != 0:
        return false
    7. q_raw ← auth_len / n
    8. if q_raw < 1 or q_raw > hsf:
        return false
    9. wots_sig ← uxmss_sig[0 : wots_sig_size]
    10. auth ← uxmss_sig[wots_sig_size : ]
    11. if q_raw < hsf:
            candidates ← {q_raw}
        else:
            candidates ← {hsf, hsf + 1}
    12. for each q in candidates:
        a. ADRS ← new_ADRS()
        b. PK.sf ← uxmss_pkFromSig(wots_sig, auth, m, PK.seed, PK.root, ADRS, q)
        c. if PK.sf == false:
            continue
        d. setLayerAddress(ADRS, 0)
            setTreeAddress(ADRS, 0)
            setTypeAndClear(ADRS, ROOT)
            expected_root ← Tw(PK.seed, ADRS, PK.sf || PK.sl)
        e. if expected_root == PK.root:
            return true
    13. return false
\end{verbatim}

\subsection{Stateless verification}
The \texttt{SHRINCS.VerifyStateless} function verifies a stateless signature. It extracts \texttt{PK.sf} from the signature, recomputes the unified digest from the FORS+C signature's \texttt{R} and \texttt{ctr}, extracts both FORS indices and hypertree indices from the same digest, then verifies bottom-up through the hypertree.
\begin{verbatim}
    Algorithm: SHRINCS.VerifyStateless(m, sig, PK)
    Input:
        m: message
        sig: stateless signature
        PK: public key
    Output:
        valid: boolean

    1.  ADRS ← new_ADRS()
    2.  PK.sf ← sig[0 : n]
    3.  offset ← n
    4.  fors_sig_size ← R_SIZE + (k - 1) * (n + a * n)
    5.  fors_sig ← sig[offset : offset + fors_sig_size]
    6.  offset ← offset + fors_sig_size
    7.  R ← fors_sig[0 : R_SIZE]
    8.  setTypeAndClear(ADRS, SL_H_MSG)
    9.  digest ← H_msg_fors(ADRS, R, PK.seed, PK.root, m)
    10. indices ← fors_msg_to_indices(digest)
    11. if indices[k-1] != 0:
            return false
    12. (tree_idx, leaf_idx) ← parse_idx(digest)
    13. setLayerAddress(ADRS, 0)
    14. setTreeAddress(ADRS, tree_idx[0] * 2^h_prime + leaf_idx[0])
    15. fors_pk ← fors_pkFromSig(fors_sig, indices, PK.seed, ADRS)
    16. msg ← fors_pk
    17. xmss_sig_size ← R_SIZE + c + l * n + h_prime * n
    18. for layer from 0 to d - 1:
        a. xmss_sig ← sig[offset : offset + xmss_sig_size]
        b. offset ← offset + xmss_sig_size
        c. wots_sig ← xmss_sig[0 : R_SIZE + c + l * n]
        d. auth ← xmss_sig[R_SIZE + c + l * n : R_SIZE + c + l * n + h' * n]
        e. setLayerAddress(ADRS, layer)
        f. setTreeAddress(ADRS, tree_idx[layer])
        g. msg ← xmss_pkFromSig(wots_sig, auth, msg, PK.seed, PK.root, ADRS, h_prime, leaf_idx[layer])
        h. if msg == false:
               return false
    19. PK.sl ← msg
    20. setLayerAddress(ADRS, 0)
    21. setTreeAddress(ADRS, 0)
    22. setTypeAndClear(ADRS, ROOT)
    23. expected_root ← Tw(PK.seed, ADRS, PK.sf || PK.sl)
    24. return (expected_root == PK.root)
\end{verbatim}

\subsection{Unified Verification}
The \texttt{SHRINCS.Verify} function provides a single entry point for verifying either stateful or stateless signatures. It distinguishes between the two types based on signature length: stateful signatures have a maximum size of \texttt{max\_sf\_size = n + R\_SIZE + c + l*n + hsf*n} bytes (both \texttt{q = hsf} and \texttt{q = hsf + 1} produce hsf auth-path nodes), while stateless signatures are strictly larger (see Section~8.7 for the derivation that this bound always holds).
\begin{verbatim}
    Algorithm: SHRINCS.Verify(m, sig, PK)
    Input:
        m: message
        sig: signature
        PK: public key
    Output:
        valid: boolean

    1. max_sf_size ← n + R_SIZE + c + l * n + hsf * n
    2. if len(sig) <= max_sf_size:
            return SHRINCS.VerifyStateful(m, sig, PK)
        else:
            return SHRINCS.VerifyStateless(m, sig, PK)
\end{verbatim}